\clearpage\SourceLeaf{293}

\begin{center}
\textbf{REMARKS ON THE GEOMETRIA SITUS NOTES.}
\end{center}

The notes collected under [I] follow one another in a
notebook, yet they were evidently written at widely differing
times. Notes [1] and [2] probably date
from the period between 1823 and 1827, while note [9] presumably
belongs to the period after 1840.

Among \textsc{Schering}'s papers were found the following explanations of the notes under
[I].

»To obtain the characters $\gamma$, $\gamma'$ [p.\,272] of any given
part of the curve, draw two straight lines from arbitrarily far away to
two points immediately beside that part of the curve and on opposite
sides of the path. For each
intersection of either of the two lines with the given path,
count $+1$ if the line passes from the outer side to the inner side, and
$-1$ in the opposite case. The results for the two lines are
the required characters $\gamma$, $\gamma'$.

If $+1$ or $-1$ is called the character of a region bounded by a single
one of the lines that return to themselves, according as the
region lies on the positive or negative side of its boundary, then
$\mathrm{Ch}\,\varepsilon a$ [p.\,274] denotes the sum of the
characters of all the regions adjacent to $a$, both those
that extend beyond $a$ and the region that is bounded by
$\varepsilon a$ and whose character is $\varepsilon$.

The absolute character $\chi$ of a point [p.\,278] is the sum of the
characters of all the regions in which the point lies.

For the count [p.\,281], one must traverse the sides of the polygon in one and
the same sense, going from $(x, y)$ to $(x', y')$
each time.«

To note [II] \textsc{Schering} assigned, probably on the basis of the statements
by \textsc{Gauss} in the note on p.\,272, the date 1844.\ Dec.\ 30.
He made the following remarks about it.

»Two points of the track that coincide and thereby form a knot
are denoted by the same letters \emph{a} and \emph{a}, \emph{b} and
\emph{b}, and\,so\,on.\ In representing a track,
these points are given in the order in which they follow one another
along the track.

The numbers

\begin{center}
1 \quad 54 \quad 15 \quad 3 \quad 10

15
\end{center}

are represented by the figures.«\hfill \textsc{Stäckel}.
